\documentclass{article}
\usepackage{amsmath,amsthm,amssymb}
\usepackage{mathtext}
\usepackage[T2A]{fontenc}
\usepackage[utf8]{inputenc}
\usepackage[english,russian]{babel}
\usepackage{graphicx}
\usepackage{hyperref}
\usepackage[margin=2.5cm]{geometry}

% ---- Итоговые числа (test macro-F1, %) ----
\newcommand{\fTfidf}{75.07}
\newcommand{\fRubert}{78.28}
\newcommand{\fRumodern}{78.57}
\newcommand{\fXlmr}{79.44}
\newcommand{\fRuroberta}{79.12}
\newcommand{\fEns}{80.14}
\newcommand{\aEns}{79.93}
\newcommand{\fEnsUni}{79.89}

\title{Анализ тональности русскоязычных отзывов о товарах: \\ ансамбль дообученных трансформеров на RuReviews}
\author{Баженов Данил Борисович}
\date{Октябрь 2026}

\begin{document}
\sloppy
\maketitle

\begin{abstract}
Мы рассматриваем задачу трёхклассовой классификации тональности (негативная / нейтральная / позитивная) русскоязычных отзывов о товарах на открытом датасете RuReviews~\cite{smetanin2019rureviews}. Мы анализируем датасет и показываем, что в стандартном протоколе оценки около 6\% тестовых отзывов дословно (после нормализации) встречаются в обучающей выборке, а 388 текстов имеют противоречивые метки; поэтому дополнительно предлагается протокол без пересечений. Мы сравниваем TF-IDF-бейзлайн, воспроизведение опубликованных подходов (RuBERT, XLM-RoBERTa-Large) и современные энкодеры (RuModernBERT, ruRoBERTa-large), а также предлагаем ансамбль с весами, подобранными на валидации. Ансамбль достигает macro-F1 = \fEns\% на тесте, что выше лучшего опубликованного результата 78.81\%~\cite{smetanin2022transformers}.
Код: \url{https://github.com/danil39russia/rureviews-sentiment}.
\end{abstract}

\section{Введение}
Анализ тональности отзывов --- одна из наиболее востребованных прикладных задач NLP: маркетплейсы используют его для мониторинга качества товаров и продавцов, ранжирования отзывов и выявления проблемных позиций. Для русского языка качество моделей заметно отстаёт от английского из-за меньшего числа размеченных корпусов~\cite{kotelnikov2021landscape}.

Особенно сложен класс <<нейтральный>>: он лежит между двумя полюсами и в отзывах часто выражается смешанными оценками (<<ткань хорошая, но размер маломерит>>). Именно на нём ошибаются все модели в работе~\cite{smetanin2021deep}.

Наш вклад:
\begin{itemize}
    \item анализ качества датасета RuReviews: дубликаты, противоречивые метки и утечка между train и test в стандартном протоколе, а также альтернативный протокол без пересечений;
    \item единый воспроизводимый пайплайн сравнения: TF-IDF, RuBERT, XLM-R-Large, RuModernBERT, ruRoBERTa-large на одном и том же разбиении;
    \item ансамбль с весами, обученными на валидации, который превосходит лучший опубликованный результат на RuReviews.
\end{itemize}

\subsection{Команда}
\textbf{Баженов Данил Борисович} --- подготовка данных, реализация и проведение всех экспериментов, написание отчёта.

\section{Обзор литературы}
\label{sec:related}
\textbf{Датасеты.} Обзор русскоязычных корпусов тональности приведён в~\cite{kotelnikov2021landscape}. Основные открытые наборы: SentiRuEval-2016 (твиты о банках и телеком-компаниях)~\cite{loukachevitch2016sentirueval}, RuSentiment (посты ВКонтакте)~\cite{rogers2018rusentiment}, RuSentiTweet~\cite{smetanin2022rusentitweet}, а также RuReviews~\cite{smetanin2019rureviews} --- отзывы о товарах категории <<Женская одежда и аксессуары>>.

\textbf{Подходы до трансформеров.} Авторы RuReviews использовали свёрточную сеть~\cite{kim2014cnn} поверх эмбеддингов Word2Vec, обученных на отзывах, с препроцессингом (нижний регистр, теги для удлинённых и капслок-слов, замена эмотиконов). Результат --- macro-F1 = 75.45\%.

\textbf{Дообучение трансформеров.} В~\cite{smetanin2021deep} дообучены Multilingual BERT~\cite{devlin2019bert}, RuBERT~\cite{kuratov2019rubert} и две версии Multilingual USE на семи русскоязычных датасетах. Для датасетов без готового разбиения использовался случайный сплит 70/15/15. На RuReviews лучший результат показал RuBERT: macro-F1 = 77.44\%. В продолжении этой работы~\cite{smetanin2022transformers} сравнены 12 предобученных трансформеров; лучший на RuReviews --- XLM-RoBERTa-Large~\cite{conneau2020xlmr} с F1 = 78.81\%, далее LaBSE (78.47) и XLM-R-Base (78.28). Этот результат мы считаем текущим SotA. В~\cite{golubev2021transfer} показано, что дополнительный трансфер с автоматически размеченных данных улучшает результаты на RuSentiment и SentiRuEval.

\textbf{Новые модели для русского языка.} После 2022 года появились ruRoBERTa-large~\cite{zmitrovich2023family}, RuModernBERT (адаптация ModernBERT~\cite{warner2024modernbert}) и многоязычные эмбеддеры E5~\cite{wang2024me5}. В бенчмарке ruMTEB~\cite{snegirev2025rumteb} есть задача RuReviewsClassification, но там модели оцениваются как замороженные эмбеддеры с логистической регрессией, поэтому числа не сравнимы с дообучением.

\textbf{LLM.} Большие языковые модели в zero-shot режиме уступают дообученным моделям на задачах тональности с тонкими границами классов~\cite{zhang2024sentiment}; альтернативой является дообучение LLM с классификационной головой через LoRA/QLoRA~\cite{hu2022lora,dettmers2023qlora,li2023labelsupervised}.

\section{Описание модели}
\label{sec:model}
Каждая модель-участник $m$ --- трансформер-энкодер с линейным классификатором поверх представления первого токена, дообученный целиком с кросс-энтропией. На выходе модель даёт логиты $z^{(m)}(x)\in\mathbb{R}^3$.

\textbf{Ансамбль.} Итоговое распределение --- взвешенное среднее логарифмов вероятностей участников:
$$\log p(y\mid x) \propto \sum_{m=1}^{M} w_m \log \mathrm{softmax}\big(z^{(m)}(x)\big)_y, \qquad w_m\ge 0,\ \sum_m w_m = 1.$$
Веса $w$ подбираются минимизацией log-loss на \emph{валидационной} выборке (метод Нелдера--Мида, параметризация через softmax). Тестовая выборка используется один раз для финальной оценки. Предсказание --- $\arg\max_y p(y\mid x)$.

Мотивация: модели с разными токенизаторами и корпусами предобучения (многоязычный XLM-R, русскоязычные ruRoBERTa и RuBERT, RuModernBERT с современной архитектурой, а также символьные n-граммы TF-IDF) ошибаются по-разному, и усреднение снижает дисперсию~\cite{dietterich2000ensemble}. Это особенно важно для шумных автоматических меток RuReviews.

Участники ансамбля:
\begin{itemize}
    \item \textbf{XLM-RoBERTa-Large} (560M параметров) --- воспроизведение текущего SotA;
    \item \textbf{ruRoBERTa-large} (355M) --- RoBERTa, предобученная на русскоязычных текстах~\cite{zmitrovich2023family};
    \item \textbf{RuModernBERT-base} (150M) --- ModernBERT для русского языка;
    \item \textbf{RuBERT-base} (180M) --- воспроизведение~\cite{smetanin2021deep};
    \item \textbf{TF-IDF + логистическая регрессия} --- словные (1--2) и символьные (2--5) n-граммы.
\end{itemize}

\section{Датасет}
RuReviews~\cite{smetanin2019rureviews} содержит отзывы о товарах категории <<Женская одежда и аксессуары>> с крупного российского маркетплейса. Метки получены автоматически из пользовательских оценок: 1--2 звезды --- негатив, 3 --- нейтрально, 4--5 --- позитив. Исходные 821k отзывов сбалансированы андерсемплингом до 30\,000 на класс. Датасет распространяется по лицензии Apache-2.0 и доступен по ссылке\footnote{\url{https://github.com/sismetanin/rureviews}}; скрипт загрузки --- \texttt{scripts/download\_data.sh}.

\textbf{Анализ качества данных.} Мы обнаружили (Табл.~\ref{tab:quality}): 2\,679 точных дубликатов; 5\,068 дубликатов после нормализации (нижний регистр, \textit{ё}$\to$\textit{е}, удаление пунктуации); 388 текстов, встречающихся с \emph{разными} метками, хотя авторы заявляют, что конфликты разрешены голосованием. Типичные примеры --- короткие шаблонные отзывы: <<Товар не пришёл>> (58 раз, метки negative и neutral), <<Товар не пришёл, деньги вернули>> и даже отзыв из одной точки <<.>> (46 раз) встречаются со всеми тремя метками.

\begin{table}[tbh!]
\begin{center}
\begin{tabular}{|l|r|}
\hline
Всего отзывов & 90\,000 \\
Точных дубликатов & 2\,679 \\
Дубликатов после нормализации & 5\,068 \\
Текстов с противоречивыми метками & 388 \\
Средняя длина, символов / слов & 130.1 / 20.0 \\
Медианная длина, символов & 91 \\
\hline
\end{tabular}
\caption{Статистика и качество данных RuReviews.}
\label{tab:quality}
\end{center}
\end{table}

\textbf{Протоколы разбиения.} Официального разбиения нет. Мы используем два протокола (Табл.~\ref{tab:splits}):
\begin{enumerate}
    \item \textbf{random} --- стратифицированный случайный сплит 70/15/15, как в~\cite{smetanin2021deep}; используется для сравнения с опубликованными результатами. При этом 830 тестовых отзывов (6.1\%) после нормализации встречаются в train, то есть часть теста модель фактически видела~\cite{elangovan2021memorization};
    \item \textbf{dedup} --- те же пропорции, но разбиваются уникальные нормализованные тексты, так что пересечений между частями нет.
\end{enumerate}

\begin{table}[tbh!]
\begin{center}
\begin{tabular}[t]{|l|ccc|c|}
\hline
 & Train & Valid & Test & Тест $\cap$ Train \\
\hline
random & 63\,000 & 13\,500 & 13\,500 & 830 \\
dedup & 62\,435 & 13\,977 & 13\,588 & 0 \\
\hline
\end{tabular}
\caption{Размеры частей для двух протоколов. Классы сбалансированы (по $\approx$1/3) во всех частях.}
\label{tab:splits}
\end{center}
\end{table}

\section{Эксперименты}
\subsection{Метрики}
Основная метрика --- macro-F1, как в~\cite{smetanin2019rureviews,smetanin2021deep}:
$$F1_c = \frac{2P_cR_c}{P_c+R_c},\qquad \text{macro-F1} = \frac{1}{3}\sum_{c\in\{neg,neu,pos\}} F1_c,$$
где $P_c$ и $R_c$ --- точность и полнота класса $c$. Дополнительно приводятся macro-precision, macro-recall, accuracy и F1 по классам.

\subsection{Постановка эксперимента}
Все модели обучаются на train, лучшая эпоха выбирается по macro-F1 на valid, результат приводится на test. Гиперпараметры: AdamW, линейный warmup 10\%, weight decay 0.01, максимальная длина 256 токенов, эффективный батч 32, bf16. Base-модели: 3 эпохи, lr $3\cdot10^{-5}$ (RuBERT) и $5\cdot10^{-5}$ (RuModernBERT). Large-модели: 8-битный AdamW, батч $8\times4$; XLM-R-Large --- 3 эпохи, lr $10^{-5}$; ruRoBERTa-large --- 1 эпоха, lr $2\cdot10^{-5}$. TF-IDF: до 300k словных и 300k символьных признаков, логистическая регрессия с $C=4$. Seed 42 для сплита и обучения. Оборудование --- одна NVIDIA RTX 5080 (16 ГБ).

\subsection{Бейзлайны}
\begin{itemize}
    \item TF-IDF + логистическая регрессия --- простой бейзлайн;
    \item CNN + Word2Vec~\cite{smetanin2019rureviews}, RuBERT~\cite{smetanin2021deep}, XLM-R-Large~\cite{smetanin2022transformers} --- опубликованные результаты. RuBERT и XLM-R-Large дополнительно воспроизведены на нашем разбиении, поскольку seed исходного сплита не опубликован.
\end{itemize}

\section{Результаты}
Сравнение с опубликованными работами приведено в Табл.~\ref{tab:results}, подробные метрики --- в Табл.~\ref{tab:detailed}.

\begin{table}[!tbh]
\centering
\begin{tabular}{|l|c|}
\hline
\textbf{Модель} & \textbf{Test macro-F1, \%} \\
\hline
\multicolumn{2}{|l|}{\textit{Опубликованные результаты (сплит 70/15/15, seed неизвестен)}} \\
CNN + Word2Vec~\cite{smetanin2019rureviews} & 75.45 \\
RuBERT~\cite{smetanin2021deep} & 77.44 \\
XLM-RoBERTa-Large~\cite{smetanin2022transformers} & 78.81 \\
\hline
\multicolumn{2}{|l|}{\textit{Наши запуски (протокол random)}} \\
TF-IDF + LogReg & \fTfidf \\
RuBERT-base (воспр.) & \fRubert \\
RuModernBERT-base & \fRumodern \\
XLM-RoBERTa-Large (воспр.) & \fXlmr \\
ruRoBERTa-large & \fRuroberta \\
Ансамбль, равные веса (абляция) & \fEnsUni \\
\textbf{Ансамбль, веса по валидации (наш)} & \textbf{\fEns} \\
\hline
\end{tabular}
\caption{Сравнение с предыдущими работами на RuReviews.}
\label{tab:results}
\end{table}


\begin{table}[!tbh]
\centering
\begin{tabular}{|l|ccc|ccc|}
\hline
Модель & P & R & Acc & F1$_{neg}$ & F1$_{neu}$ & F1$_{pos}$ \\
\hline
TF-IDF + LogReg & 75.17 & 74.99 & 74.99 & 72.18 & 65.15 & 87.87 \\
RuBERT-base & 78.90 & 78.00 & 78.00 & 75.00 & 69.92 & 89.92 \\
RuModernBERT-base & 79.41 & 78.27 & 78.27 & 74.68 & 70.67 & 90.37 \\
ruRoBERTa-large & 79.28 & 78.99 & 78.99 & 76.63 & 69.92 & 90.81 \\
XLM-RoBERTa-Large & 79.93 & 79.25 & 79.25 & 75.94 & 71.51 & 90.88 \\
Ансамбль, равные веса & 80.36 & 79.68 & 79.68 & 76.60 & 71.86 & 91.22 \\
\textbf{Ансамбль (наш)} & \textbf{80.59} & \textbf{79.93} & \textbf{79.93} & \textbf{76.79} & \textbf{72.22} & \textbf{91.40} \\
\hline
\end{tabular}
\caption{Подробные метрики на тесте (протокол random, \%): macro-precision, macro-recall, accuracy и F1 по классам. Веса ансамбля, подобранные на валидации: ruRoBERTa-large 0.33, XLM-R-Large 0.32, RuModernBERT 0.19, TF-IDF 0.11, RuBERT 0.05.}
\label{tab:detailed}
\end{table}

\textbf{Интерпретация.}
\begin{itemize}
    \item Простой TF-IDF-бейзлайн (\fTfidf) почти достигает результата CNN из оригинальной работы (75.45): датасет во многом решается лексическими признаками.
    \item Воспроизведённый RuBERT на нашем сплите (\fRubert) немного выше опубликованного (77.44). Это согласуется с разбросом от выбора сплита и с утечкой дубликатов в протоколе random.
    \item Воспроизведённый XLM-R-Large (\fXlmr) также выше опубликованного значения (78.81). Поэтому для честного сравнения ансамбль сопоставляется не только с опубликованными числами, но и с сильнейшей одиночной моделью на \emph{том же} разбиении.
    \item Ансамбль превосходит и лучший опубликованный результат (78.81), и все одиночные модели на нашем разбиении, accuracy = \aEns\%. Ансамбль с равными весами (\fEnsUni) хуже ансамбля с весами, подобранными на валидации: вклад участников неравноценен.
    \item Как и в~\cite{smetanin2021deep}, основная доля ошибок связана с нейтральным классом: F1 для него (72.2 у ансамбля) заметно ниже, чем для негатива (76.8) и позитива (91.4). Ансамбль улучшает F1 по всем трём классам. Часть ошибок объясняется шумом меток (Табл.~\ref{tab:quality}): одинаковые тексты встречаются с разными оценками, что задаёт потолок качества для любой модели.
\end{itemize}

\textbf{Ограничения.} Все результаты получены с одним seed; из-за ограничений по времени не проведены запуски с несколькими seed, оценка на протоколе dedup, а также эксперименты с LLM (Qwen2.5-7B: zero-/few-shot и QLoRA-классификатор). Код для них есть в репозитории (\texttt{src/zero\_shot\_llm.py}, \texttt{src/train\_llm\_lora.py}).

\section{Заключение}
Мы проанализировали датасет RuReviews и нашли в нём дубликаты, противоречивые метки и утечку между train и test в стандартном протоколе оценки; предложен протокол без пересечений. В едином воспроизводимом пайплайне сравнены TF-IDF, RuBERT, XLM-R-Large, RuModernBERT и ruRoBERTa-large. Предложенный ансамбль с весами, подобранными на валидации, достигает macro-F1 = \fEns\% и превосходит лучший опубликованный результат на RuReviews (78.81\%). Дальнейшая работа: оценка на протоколе dedup, несколько seed и включение дообученных LLM в ансамбль.

\bibliographystyle{apalike}
\bibliography{lit}
\end{document}
